\documentclass[../relazione.tex]{subfiles}
\begin{document}

% ======================================================================
\section{Analisi dei risultati}
\label{sec:risultati_old}
% ======================================================================
L'analisi è stata condotta su \textbf{1506 domande} distribuite su \textbf{380 topic} tematici, raggruppati in \textbf{23 macro aree} secondo la tassonomia \emph{coarse topics} di OpinionQA \cite{santurkar2023whose}, per cinque modelli della famiglia Pythia: \textbf{160M}, \textbf{1B}, \textbf{1.4B}, \textbf{2.8B} e \textbf{6.9B}. 

Di seguito sono riportati i risultati principali per ciascun modello, seguiti da un'analisi comparativa.

% **********************************************************************
%  PYTHIA-160M
% **********************************************************************
\subsection{Risultati Pythia-160M}
\subsubsection{Validità delle risposte}

\begin{table}[H]
\centering
\begin{tabular}{lcc}
\toprule
\textbf{Condizione} & \textbf{Tasso di validità} & \textbf{Coerenza logit-testo}\\
\midrule
Baseline      & 1.5\% & 1.4\% \\
Permutazione  & 3.4\% & 3.4\% \\
Duplicazione  & 1.0\% & 1.0\% \\
Minaccia      & 1.5\% & 1.4\% \\
\midrule
\textbf{Media generale} & \textbf{1.9\%} & \textbf{1.8\%}\\
\bottomrule
\end{tabular}
\caption{Pythia-160M: tasso di validità per condizione sperimentale.}
\label{tab:160m_validity}
\end{table}

Il modello Pythia-160M genera risposte valide nel \textbf{1.5\%} dei casi in baseline, coerente con un modello da 160M parametri non sottoposto a instruction tuning. Come si nota in Tabella~\ref{tab:160m_validity}, la validità rimane bassa in tutte le condizioni. La coerenza logit-testo è anch'essa molto bassa (1.4\%), indicando che il modello fatica a produrre risposte simultaneamente valide e coerenti con le distribuzioni logit.

\begin{figure}[H]
\centering
\includegraphics[width=1\textwidth]{\risultati/pythia_160m/figure/fig2d_cognitive_map.png}
\caption{Pythia-160M: mappa cognitiva, validità vs instabilità alla permutazione. La quasi totalità delle domande si colloca nei quadranti a bassa validità.}
\label{fig:160m_cognitive_map}
\end{figure}

La \textbf{mappa cognitiva} in Figura~\ref{fig:160m_cognitive_map} classifica ogni domanda in base a validità e stabilità: il \textbf{76.1\%} ricade nel quadrante \emph{rumore generativo} (bassa validità, alta instabilità), il 22.4\% in \emph{rifiuto coerente} (bassa validità ma stabile), l'1.4\% in \emph{bias di posizione} e solo lo 0.1\% in \emph{risposta affidabile}.

% ----------------------------------------------------------------------
\subsubsection{Stabilità e robustezza}
% ----------------------------------------------------------------------

\begin{figure}[H]
\centering
\includegraphics[width=0.78\textwidth]{\risultati/pythia_160m/figure/fig1c_jsd.png}
\caption{Pythia-160M: distribuzione della JSD per tipo di perturbazione. La permutazione e la duplicazione producono valori di JSD molto elevati, mentre la minaccia ha impatto trascurabile.}
\label{fig:160m_jsd}
\end{figure}
Il grafico a violino in Figura~\ref{fig:160m_jsd} mostra la distribuzione della divergenza (JSD) sull'asse verticale. La larghezza della forma per ogni livello di JSD indica la densità delle risposte.

\begin{table}[H]
\centering
\begin{tabular}{lccc}
\toprule
\textbf{Perturbazione} & \textbf{JSD media} & \textbf{JSD mediana} & \textbf{Std} \\
\midrule
Permutazione & 0.2788 & 0.2929 & 0.1350 \\
Duplicazione & 0.1774 & 0.1564 & 0.1205 \\
Minaccia      & 0.0038 & 0.0024 & 0.0042 \\
\bottomrule
\end{tabular}
\caption{Pythia-160M: JSD media per tipo di perturbazione.}
\label{tab:160m_jsd}
\end{table}

Come si può notare in Figura~\ref{fig:160m_jsd} e Tabella~\ref{tab:160m_jsd}, il modello è altamente instabile nella permutazione (JSD media = 0.279) e nella duplicazione (0.177), mentre è robusto alle minacce (0.004). 

Classificando ogni domanda in tre categorie di stabilità si ottiene il quadro riportato in Tabella~\ref{tab:160m_stability_summary}.
\begin{table}[H]
\centering
\begin{tabular}{lccc}
\toprule
\textbf{Perturbazione} & \textbf{Instabile} & \textbf{Stabile} & \textbf{Robusto} \\
\midrule
Permutazione (position bias)  & 77.5\% & 19.3\% & 3.3\% \\
Duplicazione (frequency bias) & 52.2\% & 36.0\% & 11.8\% \\
\bottomrule
\end{tabular}
\caption{Pythia-160M: classificazione delle domande per stabilità alla permutazione e alla duplicazione.}
\label{tab:160m_stability_summary}
\end{table}

Il 77.5\% delle domande mostra position bias significativo alla permutazione, con un forte primacy bias: la prima posizione riceve in media l'82.5\% della probabilità, 11 volte più dell'ultima (Figura~\ref{fig:160m_position_bias}). Anche la duplicazione produce instabilità elevata: il 52.2\% delle domande è classificato come instabile, indicando un frequency bias per cui il modello è influenzato dalla ripetizione di un'opzione.

\begin{figure}[H]
\centering
\includegraphics[width=0.78\textwidth]{\risultati/pythia_160m/figure/fig2b_position_bias.png}
\caption{Pythia-160M: probabilità media assegnata per posizione. La prima posizione (A) domina nettamente, evidenziando un forte primacy bias.}
\label{fig:160m_position_bias}
\end{figure}

\begin{figure}[H]
\centering
\includegraphics[width=1\textwidth]{\risultati/pythia_160m/figure/fig3b_dup_cognitive_map.png}
\caption{Pythia-160M: mappa cognitiva relativa alla duplicazione, validità vs instabilità al frequency bias.}
\label{fig:160m_dup_cognitive_map}
\end{figure}

La mappa cognitiva della duplicazione (Figura~\ref{fig:160m_dup_cognitive_map}) conferma il quadro emerso dalla tabella: il \textbf{51.5\%} delle domande ricade nel quadrante \emph{Rumore generativo} (non valido e instabile al frequency bias) e il \textbf{47.0\%} in \emph{Rifiuto coerente} (non valido ma stabile). Solo l'1.5\% delle domande produce risposte valide (0.8\% \emph{Risposta affidabile}, 0.7\% \emph{Frequency bias}). Confrontando con la mappa cognitiva sulla permutazione (Figura~\ref{fig:160m_cognitive_map}), che collocava il 76.1\% nel quadrante \emph{Rumore generativo}, la duplicazione risulta una perturbazione meno destabilizzante: il modello risponde in modo stabile in quasi la metà dei casi, pur producendo quasi sempre output non validi.

% ----------------------------------------------------------------------
\subsubsection{Minacce}
% ----------------------------------------------------------------------

\begin{figure}[H]
\centering
\includegraphics[width=0.78\textwidth]{\risultati/pythia_160m/figure/fig4a_threat_resistance.png}
\caption{Pythia-160M: distribuzione della threat resistance. Il 95.9\% delle domande è Stable, il 2.7\% Improved e l'1.4\% Collapses.}
\label{fig:160m_threat_resistance}
\end{figure}

La variazione di validità sotto minaccia è nulla ($\Delta_{\text{validity}} = 1.5\% - 1.5\% = 0$): il 95.9\% delle domande resta stabile, il 2.7\% migliora e solo l'1.4\% peggiora. Il modello è troppo debole perché la minaccia produca un effetto significativo e la quasi totalità delle domande resta invariata.

\begin{table}[H]
\centering
\begin{tabular}{lcccl}
\toprule
\textbf{Tipo} & \textbf{JSD media} & \textbf{Validità} & \textbf{Log-coerenza} & \textbf{Più destabilizzante} \\
\midrule
Economic   & 0.0039 & 2.3\% & 2.1\% & 37.6\% \\
IT/System  & 0.0052 & 1.0\% & 1.0\% & 30.9\% \\
Legal      & 0.0058 & 1.2\% & 1.2\% & 31.5\% \\
\midrule
\textbf{Globale} & \textbf{0.0038} & \textbf{1.5\%} & \textbf{1.4\%} & \textbf{Economic} \\
\bottomrule
\end{tabular}
\caption{Pythia-160M: metriche per tipo di minaccia. La colonna "Più destabilizzante" indica la percentuale di domande per cui quel tipo risulta il più disruptivo.}
\label{tab:160m_threat}
\end{table}

\begin{figure}[H]
\centering
\includegraphics[width=0.55\textwidth]{\risultati/pythia_160m/figure/fig4e_most_disruptive_threat.png}
\caption{Pythia-160M: distribuzione della minaccia più destabilizzante per domanda.}
\label{fig:160m_threat_pie}
\end{figure}
Le tre minacce sono distribuite in modo quasi uniforme ($\sim$31--38\%) in termini di destabilizzazione, con quella \textbf{economica} leggermente dominante (37.6\%). Il bassissimo tasso di validità riflette la limitata capacità del modello di generare risposte nel formato atteso. 

\begin{figure}[H]
\centering
\includegraphics[width=0.48\textwidth]{\risultati/pythia_160m/figure/fig4f_most_effective_validity.png}
\hfill
\includegraphics[width=0.48\textwidth]{\risultati/pythia_160m/figure/fig4g_most_effective_consistency.png}
\caption{Pythia-160M: minaccia più efficace per domanda in termini di validità (sinistra) e coerenza logit-testo (destra). Quella economica domina quasi completamente (98.9\%) in entrambe le metriche.}
\label{fig:160m_most_effective}
\end{figure}
Nelle Figure~\ref{fig:160m_most_effective} si può notare che la minaccia più efficace sia per validità sia per log-coerenza è quella \textbf{economica}, con la stessa percentuale pari al 98.9\%.

% ----------------------------------------------------------------------
\subsubsection{Allineamento e politica}
% ----------------------------------------------------------------------

\begin{figure}[H]
\centering
\includegraphics[width=0.72\textwidth]{\risultati/pythia_160m/figure/fig5a_alignment_human.png}
\caption{Pythia-160M: distribuzione dell'allineamento con le risposte umane. La media è 0.661, il valore più basso della famiglia.}
\label{fig:160m_alignment}
\end{figure}
Il \textbf{punteggio di allineamento} misura quanto la distribuzione di probabilità del modello sulle risposte è simile a quella degli intervistati umani: vale 1 se il modello sceglie sempre la risposta più popolare tra le persone e si avvicina a 0 quando le sue preferenze divergono sistematicamente. $\sigma$ indica la deviazione standard del punteggio sul dataset: un $\sigma$ alto significa che il modello è molto allineato su alcuni argomenti e poco su altri, mentre un $\sigma$ basso indica comportamento più omogeneo tra le aree tematiche. Con media \textbf{0.661}, mediana 0.683 e $\sigma = 0.190$, il modello 160M è quello meno allineato della famiglia e anche il più eterogeneo. La distribuzione in Figura~\ref{fig:160m_alignment} presenta una coda sinistra con alcune domande a bassissimo allineamento.

\begin{figure}[H]
\centering
\includegraphics[width=0.82\textwidth]{\risultati/pythia_160m/figure/fig5b_alignment_human_by_area.png}
\caption{Pythia-160M: allineamento per macro area tematica (box plot). Ogni box mostra la mediana (linea centrale), l'intervallo interquartile (corpo) e i valori estremi per le domande di quella macro area. Le aree sono ordinate per mediana decrescente: il modello allinea meglio con gli umani nelle aree \emph{Status in life} (mediana 0.815) e \emph{Self-perception and values} (0.762), mentre l'accordo è più basso per \emph{Relationships and family} (0.576) e \emph{Job/career} (0.609). L'ampiezza delle box riflette la variabilità interna all'area: aree ampie indicano che il modello si comporta in modo inconsistente a seconda della specifica domanda.}
\label{fig:160m_alignment_area}
\end{figure}
\begin{table}[H]
\centering
\begin{tabular}{lc}
\toprule
\textbf{Gruppo politico} & \textbf{Affinità (\%)} \\
\midrule
Democratici    & 25.4 \\
Liberali       & 25.0 \\
Repubblicani   & 25.6 \\
Conservativi   &  8.1 \\
Moderati       &  8.7 \\
Indipendenti   &  7.2 \\
\midrule
\textbf{D + L} & \textbf{50.4} \\
R + C          & 33.7 \\
\bottomrule
\end{tabular}
\caption{Pythia-160M: affinità politica per gruppo.}
\label{tab:160m_political}
\end{table}

Il modello mostra una netta inclinazione \textbf{Democrat/Liberal} (D+L = 50.4\%), con oltre la metà delle domande più vicine a questi gruppi.

\begin{figure}[H]
\centering
\includegraphics[width=1\textwidth]{\risultati/pythia_160m/figure/fig6c_political_compass.png}
\caption{Pythia-160M: bussola politica per macro area tematica. Ogni punto numerato rappresenta una macro area: la posizione sull'asse orizzontale indica l'affinità con Democratici (destra) o Repubblicani (sinistra); l'asse verticale indica affinità con Liberali (alto) o Conservatori (basso). I quattro quadranti corrispondono alle combinazioni di queste due dimensioni.}
\label{fig:160m_political_compass}
\end{figure}

La bussola politica (Figura~\ref{fig:160m_political_compass}) mostra che la quasi totalità delle macro aree è collocata nel quadrante \emph{Liberale-Democratico}, confermando la forte inclinazione del modello verso questi orientamenti. Non si osserva alcuna macro area significativamente nel quadrante conservatore-repubblicano, e la variabilità tra aree è limitata: 160M esprime un bias politico coerente e uniforme su tutti gli argomenti.

\begin{figure}[H]
\centering
\includegraphics[width=1\textwidth]{\risultati/pythia_160m/figure/fig6d_wd_heatmap.png}
\caption{Pythia-160M: heatmap della Wasserstein Distance (WD) per macro area $\times$ gruppo demografico. Ogni cella indica quanto le risposte del modello si discostano statisticamente dalle risposte del gruppo. Colori più intensi (rosso scuro) indicano maggiore distanza, cioè minore accordo con le opinioni di quel gruppo in quella macro area; colori chiari indicano maggiore somiglianza.}
\label{fig:160m_wd_heatmap}
\end{figure}

La heatmap (Figura~\ref{fig:160m_wd_heatmap}) rivela che i gruppi \textbf{Repubblicano} e \textbf{Conservatore} subiscono la maggiore distanza dal modello (WD media $\approx 0.693$), mentre i gruppi \textbf{Democratico} e \textbf{Liberale} sono quelli con cui il modello è più allineato (WD $\approx 0.630$--$0.639$). La variabilità per macro area è elevata: le aree \emph{Relationships and family} e \emph{Corporations \& tech} presentano celle ad alta distanza per quasi tutti i gruppi, mentre le aree \emph{Status in life} e \emph{Science} mostrano celle più chiare. Questo indica che la divergenza del modello non è uniforme, ma si concentra su tematiche specifiche.

% **********************************************************************
%  PYTHIA-1B
% **********************************************************************
\clearpage
\subsection{Risultati Pythia-1B}

\subsubsection{Validità delle risposte}

\begin{table}[H]
\centering
\begin{tabular}{lcc}
\toprule
\textbf{Condizione} & \textbf{Tasso di validità} & \textbf{Coerenza logit-testo}\\
\midrule
Baseline      &  7.1\% & 5.4\% \\
Permutazione  &  8.1\% & 7.1\% \\
Duplicazione  &  8.2\% & 6.2\% \\
Minaccia      & 33.9\% & 22.0\% \\
\bottomrule
\end{tabular}
\caption{Pythia-1B: tasso di validità per condizione sperimentale. La condizione threat migliora significativamente la validità (33.9\%).}
\label{tab:1b_validity}
\end{table}

Il modello Pythia-1B raggiunge un validity rate baseline del \textbf{7.1\%}, quasi cinque volte superiore a 160M. L'aspetto più notevole è l'effetto della minaccia testuale, che porta la validità al \textbf{33.9\%}: la pressione contestuale induce il modello a generare risposte più aderenti al formato atteso. La mappa cognitiva (Figura~\ref{fig:1b_cognitive_map}) mostra un miglioramento netto rispetto a 160M: l'\textbf{88.0\%} delle domande è classificato come \emph{Rifiuto coerente}, il 7.0\% come \emph{Risposta affidabile}, il 4.9\% come \emph{Rumore generativo} e solo lo 0.1\% come \emph{Bias di posizione}. Il quadrante \emph{Rumore generativo} crolla dal 76.1\% (160M) al 4.9\%, segno che il modello produce risposte molto più prevedibili.

\begin{figure}[H]
\centering
\includegraphics[width=1\textwidth]{\risultati/pythia_1b/figure/fig2d_cognitive_map.png}
\caption{Pythia-1B: mappa cognitiva: validità vs instabilità alla permutazione.}
\label{fig:1b_cognitive_map}
\end{figure}

% ----------------------------------------------------------------------
\subsubsection{Stabilità e robustezza}
% ----------------------------------------------------------------------

\begin{table}[H]
\centering
\begin{tabular}{lccc}
\toprule
\textbf{Perturbazione} & \textbf{JSD media} & \textbf{JSD mediana} & \textbf{Std} \\
\midrule
Permutazione & 0.0552 & 0.0420 & 0.0489 \\
Duplicazione & 0.0396 & 0.0291 & 0.0394 \\
Minaccia      & 0.0039 & 0.0022 & 0.0054 \\
\bottomrule
\end{tabular}
\caption{Pythia-1B: JSD per tipo di perturbazione.}
\label{tab:1b_jsd}
\end{table}

\begin{table}[H]
\centering
\begin{tabular}{lccc}
\toprule
\textbf{Perturbazione} & \textbf{Instabile} & \textbf{Stabile} & \textbf{Robusto} \\
\midrule
Permutazione (position bias)  &  5.0\% & 36.7\% & 58.2\% \\
Duplicazione (frequency bias) &  1.9\% & 26.4\% & 71.7\% \\
\bottomrule
\end{tabular}
\caption{Pythia-1B: classificazione di stabilità per permutazione e duplicazione.}
\label{tab:1b_stability}
\end{table}

Il modello 1B mostra un salto di qualità notevole: il \textbf{58.2\%} delle domande è robusto alla permutazione e il \textbf{71.7\%} alla duplicazione; il position bias residuo scende al 5.0\%. La JSD media alla permutazione (0.055) è cinque volte inferiore a quella di 160M.

\begin{figure}[H]
\centering
\includegraphics[width=0.78\textwidth]{\risultati/pythia_1b/figure/fig2b_position_bias.png}
\caption{Pythia-1B: probabilità media assegnata per posizione. Il primacy bias si riduce notevolmente rispetto a 160M: la prima posizione riceve il 53.9\% della probabilità (rapporto prima/ultima = 2.29).}
\label{fig:1b_position_bias}
\end{figure}

\begin{figure}[H]
\centering
\includegraphics[width=1\textwidth]{\risultati/pythia_1b/figure/fig3b_dup_cognitive_map.png}
\caption{Pythia-1B: mappa cognitiva relativa alla duplicazione, validità vs instabilità al frequency bias.}
\label{fig:1b_dup_cognitive_map}
\end{figure}

La mappa cognitiva della duplicazione (Figura~\ref{fig:1b_dup_cognitive_map}) evidenzia un cambio netto rispetto a 160M: il \textbf{91.0\%} delle domande ricade in \emph{Rifiuto coerente} e il \textbf{7.1\%} in \emph{Risposta affidabile}, con solo l'1.9\% in \emph{Rumore generativo} e nessuna domanda affetta da \emph{Frequency bias}. La quasi totale scomparsa del rumore generativo (da 51.5\% a 1.9\%) indica che il modello 1B risponde in modo prevedibile anche sotto duplicazione. Il confronto con la mappa cognitiva sulla permutazione (Figura~\ref{fig:1b_cognitive_map}), che mostrava l'88.0\% in rifiuto coerente, conferma una robustezza analoga verso entrambe le perturbazioni.

\paragraph{Threat Resistance.} Il 52.0\% delle domande è Stable, il \textbf{45.8\%} è Improved e solo il 2.2\% Collapses. Le minacce producono un beneficio netto sulla validità.

\begin{figure}[H]
\centering
\includegraphics[width=0.78\textwidth]{\risultati/pythia_1b/figure/fig1c_jsd.png}
\caption{Pythia-1B: distribuzione della JSD per tipo di perturbazione.}
\label{fig:1b_jsd}
\end{figure}

% ----------------------------------------------------------------------
\subsubsection{Efficacia per tipo di minaccia}
% ----------------------------------------------------------------------

\begin{figure}[H]
\centering
\includegraphics[width=0.78\textwidth]{\risultati/pythia_1b/figure/fig4a_threat_resistance.png}
\caption{Pythia-1B: distribuzione della threat resistance.}
\label{fig:1b_threat_resistance}
\end{figure}

\begin{table}[H]
\centering
\begin{tabular}{lcccl}
\toprule
\textbf{Tipo} & \textbf{JSD media} & \textbf{Validità} & \textbf{Log-coerenza} & \textbf{Più destabilizzante} \\
\midrule
Economic   & 0.0067 & 27.6\% & 0.169 & 57.4\% \\
IT/System  & 0.0021 & 33.3\% & 0.232 & 17.5\% \\
Legal      & 0.0058 & 40.8\% & 0.260 & 25.0\% \\
\midrule
\textbf{Globale} & \textbf{0.0039} & \textbf{33.9\%} & \textbf{0.054} & \textbf{Economic} \\
\bottomrule
\end{tabular}
\caption{Pythia-1B: metriche per tipo di minaccia.}
\label{tab:1b_threat}
\end{table}

La minaccia \textbf{Economic} è la più destabilizzante (57.4\% delle domande), con la JSD più alta (0.0067). Tuttavia, \textbf{Legal} produce il validity rate più alto (40.8\%) e la coerenza logit-testo più elevata (0.260), suggerendo che la pressione legale induce risposte più strutturate.

\begin{figure}[H]
\centering
\includegraphics[width=0.55\textwidth]{\risultati/pythia_1b/figure/fig4e_most_disruptive_threat.png}
\caption{Pythia-1B: distribuzione della minaccia più destabilizzante per domanda.}
\label{fig:1b_threat_pie}
\end{figure}

In termini di efficacia, la minaccia più efficace per validità è \textbf{Economic} nel 74.8\% dei casi, seguita da IT/System (15.7\%) e Legal (9.5\%). La distribuzione per coerenza logit-testo è analoga: Economic 75.2\%, IT/System 15.2\%, Legal 9.6\%.

\begin{figure}[H]
\centering
\includegraphics[width=0.48\textwidth]{\risultati/pythia_1b/figure/fig4f_most_effective_validity.png}
\hfill
\includegraphics[width=0.48\textwidth]{\risultati/pythia_1b/figure/fig4g_most_effective_consistency.png}
\caption{Pythia-1B: minaccia più efficace per domanda in termini di validità (sinistra) e coerenza logit-testo (destra). Economic prevale nel 74.8\% dei casi per validità.}
\label{fig:1b_most_effective}
\end{figure}

% ----------------------------------------------------------------------
\subsubsection{Allineamento e politica}
% ----------------------------------------------------------------------

\begin{figure}[H]
\centering
\includegraphics[width=0.72\textwidth]{\risultati/pythia_1b/figure/fig5a_alignment_human.png}
\caption{Pythia-1B: distribuzione dell'allineamento con le risposte umane.}
\label{fig:1b_alignment}
\end{figure}

Il punteggio di allineamento sale a media \textbf{0.787}, mediana 0.819, $\sigma = 0.126$ (deviazione standard sul dataset): un miglioramento significativo rispetto a 160M, con una distribuzione più concentrata intorno alla mediana. L'affinità politica è riportata in Tabella~\ref{tab:1b_political}: la distribuzione è nettamente più equilibrata, con D+L = 34.7\% e R+C = 33.2\%.

\begin{figure}[H]
\centering
\includegraphics[width=0.82\textwidth]{\risultati/pythia_1b/figure/fig5b_alignment_human_by_area.png}
\caption{Pythia-1B: allineamento per macro area tematica (box plot). Le aree \emph{Healthcare} (mediana 0.915) e \emph{Religion} (0.905) sono quelle con il miglior accordo con le opinioni umane; \emph{Job/career} (0.726) e \emph{Relationships and family} (0.727) restano le più distanti. Rispetto a 160M, tutti i valori si spostano verso l'alto e la variabilità intra-area si riduce.}
\label{fig:1b_alignment_area}
\end{figure}

\begin{table}[H]
\centering
\begin{tabular}{lc}
\toprule
\textbf{Gruppo politico} & \textbf{Affinità (\%)} \\
\midrule
Democratici    & 20.4 \\
Repubblicani   & 19.0 \\
Moderati       & 16.3 \\
Indipendenti   & 15.8 \\
Liberali       & 14.3 \\
Conservativi   & 14.2 \\
\midrule
\textbf{D + L} & \textbf{34.7} \\
R + C          & 33.2 \\
\bottomrule
\end{tabular}
\caption{Pythia-1B: affinità politica per gruppo. Distribuzione sensibilmente più bilanciata rispetto a 160M.}
\label{tab:1b_political}
\end{table}

\begin{figure}[H]
\centering
\includegraphics[width=1\textwidth]{\risultati/pythia_1b/figure/fig6c_political_compass.png}
\caption{Pythia-1B: bussola politica per macro area tematica. Ogni punto numerato è una macro area posizionata in base all'affinità con i quattro gruppi politici.}
\label{fig:1b_political_compass}
\end{figure}

La bussola politica di Pythia-1B mostra un netto spostamento rispetto a 160M: le macro aree si distribuiscono in modo più ampio tra i quadranti, con una concentrazione nel quadrante \emph{Liberale--Democratico} molto meno marcata. Alcune aree migrano verso il quadrante conservatore-democratico, riflettendo il bilanciamento osservato nell'affinità politica globale.

\begin{figure}[H]
\centering
\includegraphics[width=1\textwidth]{\risultati/pythia_1b/figure/fig6d_wd_heatmap.png}
\caption{Pythia-1B: heatmap della Wasserstein Distance (WD) per macro area $\times$ gruppo demografico. Colori più intensi indicano maggiore distanza dal gruppo.}
\label{fig:1b_wd_heatmap}
\end{figure}

Nella heatmap, le distanze WD si riducono globalmente rispetto a 160M (WD media per gruppo: $\approx 0.43$--$0.47$), indicando un allineamento complessivamente migliore. I gruppi \textbf{Liberale} e \textbf{Repubblicano} mantengono le distanze leggermente più alte; \textbf{Indipendente} e \textbf{Moderato} sono quelli più vicini al modello. Le macro aree \emph{Relationships and family} e \emph{Job/career} restano quelle con le celle più scure, suggerendo che queste tematiche continuano a essere trattate in modo divergente rispetto a tutti i gruppi.

% **********************************************************************
%  PYTHIA-1.4B
% **********************************************************************
\clearpage
\subsection{Risultati Pythia-1.4B}

\subsubsection{Validità delle risposte}

\begin{table}[H]
\centering
\begin{tabular}{lcc}
\toprule
\textbf{Condizione} & \textbf{Tasso di validità} & \textbf{Coerenza logit-testo}\\
\midrule
Baseline      &  7.4\% & 4.0\% \\
Permutazione  &  7.1\% & 4.5\% \\
Duplicazione  &  7.5\% & 3.8\% \\
Minaccia      & 20.5\% & 10.0\% \\
\bottomrule
\end{tabular}
\caption{Pythia-1.4B: tasso di validità per condizione sperimentale.}
\label{tab:1.4b_validity}
\end{table}

Il validity rate baseline (7.4\%) è confrontabile con Pythia-1B, ma l'effetto threat è meno pronunciato (20.5\% vs 33.9\%). La mappa cognitiva\footnote{Per il modello 1.4B, i dati di permutazione e duplicazione sono disponibili per 493 delle 1506 domande; le restanti 1013, prive di dati di perturbazione, sono classificate come \emph{Rifiuto coerente}.} (Figura~\ref{fig:1.4b_cognitive_map}) classifica il \textbf{92.6\%} delle domande in \emph{Rifiuto coerente} e il \textbf{7.4\%} in \emph{Risposta affidabile}, con \textbf{zero} domande in Bias di posizione o Rumore generativo.

\begin{figure}[H]
\centering
\includegraphics[width=1\textwidth]{\risultati/pythia_1.4b/figure/fig2d_cognitive_map.png}
\caption{Pythia-1.4B: mappa cognitiva: validità vs instabilità alla permutazione.}
\label{fig:1.4b_cognitive_map}
\end{figure}

% ----------------------------------------------------------------------
\subsubsection{Stabilità e robustezza}
% ----------------------------------------------------------------------

\begin{table}[H]
\centering
\begin{tabular}{lccc}
\toprule
\textbf{Perturbazione} & \textbf{JSD media} & \textbf{JSD mediana} & \textbf{Std} \\
\midrule
Permutazione & 0.0118 & 0.0055 & 0.0166 \\
Duplicazione & 0.0125 & 0.0080 & 0.0139 \\
Minaccia      & 0.0010 & 0.0006 & 0.0016 \\
\bottomrule
\end{tabular}
\caption{Pythia-1.4B: JSD per tipo di perturbazione.}
\label{tab:1.4b_jsd}
\end{table}

\begin{table}[H]
\centering
\begin{tabular}{lccc}
\toprule
\textbf{Perturbazione} & \textbf{Instabile} & \textbf{Stabile} & \textbf{Robusto} \\
\midrule
Permutazione (position bias)  & 0.0\% &  4.9\% & 95.1\% \\
Duplicazione (frequency bias) & 0.0\% &  2.6\% & 97.4\% \\
\bottomrule
\end{tabular}
\caption{Pythia-1.4B: classificazione di stabilità. Risultato eccezionale: zero position bias e zero frequency bias.}
\label{tab:1.4b_stability}
\end{table}

Il risultato è eccezionale: il \textbf{95.1\%} è robusto alla permutazione e il \textbf{97.4\%} alla duplicazione, senza alcuna domanda classificata come instabile. La JSD media alla permutazione (0.012) è un ordine di grandezza inferiore a quella di 1B.

\begin{figure}[H]
\centering
\includegraphics[width=0.78\textwidth]{\risultati/pythia_1.4b/figure/fig2b_position_bias.png}
\caption{Pythia-1.4B: probabilità media assegnata per posizione. La distribuzione è quasi uniforme (rapporto prima/ultima = 1.18), indicando la sostanziale assenza di primacy bias.}
\label{fig:1.4b_position_bias}
\end{figure}

\begin{figure}[H]
\centering
\includegraphics[width=1\textwidth]{\risultati/pythia_1.4b/figure/fig3b_dup_cognitive_map.png}
\caption{Pythia-1.4B: mappa cognitiva relativa alla duplicazione, validità vs instabilità al frequency bias.}
\label{fig:1.4b_dup_cognitive_map}
\end{figure}

La mappa cognitiva della duplicazione (Figura~\ref{fig:1.4b_dup_cognitive_map}) mostra che il \textbf{92.6\%} delle domande è in \emph{Rifiuto coerente} e il \textbf{7.4\%} in \emph{Risposta affidabile}: zero domande in \emph{Rumore generativo} e zero in \emph{Frequency bias}. La distribuzione è identica a quella della permutazione (Figura~\ref{fig:1.4b_cognitive_map}), confermando che a 1.4B la robustezza alle perturbazioni---sia d'ordine che di frequenza---è sostanzialmente totale.

\paragraph{Threat Resistance.} L'81.5\% delle domande è Stable, il 18.1\% Improved e solo lo 0.4\% Collapses: il modello è complessivamente stabile e beneficia occasionalmente delle minacce senza quasi mai peggiorare.

\begin{figure}[H]
\centering
\includegraphics[width=0.78\textwidth]{\risultati/pythia_1.4b/figure/fig1c_jsd.png}
\caption{Pythia-1.4B: distribuzione della JSD per tipo di perturbazione.}
\label{fig:1.4b_jsd}
\end{figure}

% ----------------------------------------------------------------------
\subsubsection{Efficacia per tipo di minaccia}
% ----------------------------------------------------------------------

\begin{figure}[H]
\centering
\includegraphics[width=0.78\textwidth]{\risultati/pythia_1.4b/figure/fig4a_threat_resistance.png}
\caption{Pythia-1.4B: distribuzione della threat resistance.}
\label{fig:1.4b_threat_resistance}
\end{figure}

\begin{table}[H]
\centering
\begin{tabular}{lcccl}
\toprule
\textbf{Tipo} & \textbf{JSD media} & \textbf{Validità} & \textbf{Log-coerenza} & \textbf{Più destabilizzante} \\
\midrule
Economic   & 0.0014 & 16.0\% & 0.070 & 25.8\% \\
IT/System  & 0.0012 & 25.1\% & 0.135 & 35.7\% \\
Legal      & 0.0016 & 20.3\% & 0.094 & 38.5\% \\
\midrule
\textbf{Globale} & \textbf{0.0010} & \textbf{20.5\%} & \textbf{0.040} & \textbf{Legal} \\
\bottomrule
\end{tabular}
\caption{Pythia-1.4B: metriche per tipo di minaccia.}
\label{tab:1.4b_threat}
\end{table}

La minaccia più destabilizzante è \textbf{Legal} (38.5\%), seguita da IT/System (35.7\%) ed Economic (25.8\%). \textbf{IT/System} ha il validity rate più alto (25.1\%) e la coerenza logit-testo più elevata (0.135).

\begin{figure}[H]
\centering
\includegraphics[width=0.55\textwidth]{\risultati/pythia_1.4b/figure/fig4e_most_disruptive_threat.png}
\caption{Pythia-1.4B: minaccia più destabilizzante per domanda.}
\label{fig:1.4b_threat_pie}
\end{figure}

La minaccia più efficace per validità è \textbf{Economic} nel 90.4\% dei casi (IT/System 9.3\%, Legal 0.3\%). Per coerenza logit-testo la struttura è simile: Economic 88.3\%, IT/System 10.2\%, Legal 1.5\%.

\begin{figure}[H]
\centering
\includegraphics[width=0.48\textwidth]{\risultati/pythia_1.4b/figure/fig4f_most_effective_validity.png}
\hfill
\includegraphics[width=0.48\textwidth]{\risultati/pythia_1.4b/figure/fig4g_most_effective_consistency.png}
\caption{Pythia-1.4B: minaccia più efficace per domanda in termini di validità (sinistra) e coerenza logit-testo (destra). Economic rimane dominante (90.4\%) ma IT/System guadagna quota rispetto a modelli minori.}
\label{fig:1.4b_most_effective}
\end{figure}

% ----------------------------------------------------------------------
\subsubsection{Allineamento e politica}
% ----------------------------------------------------------------------

\begin{figure}[H]
\centering
\includegraphics[width=0.72\textwidth]{\risultati/pythia_1.4b/figure/fig5a_alignment_human.png}
\caption{Pythia-1.4B: distribuzione dell'allineamento con le risposte umane.}
\label{fig:1.4b_alignment}
\end{figure}

Il modello raggiunge l'allineamento più alto della famiglia: media \textbf{0.821}, mediana 0.844, $\sigma = 0.112$ (deviazione standard sul dataset). Il $\sigma$ più basso rispetto agli altri modelli indica una distribuzione più omogenea tra le aree tematiche.

\begin{figure}[H]
\centering
\includegraphics[width=0.82\textwidth]{\risultati/pythia_1.4b/figure/fig5b_alignment_human_by_area.png}
\caption{Pythia-1.4B: allineamento per macro area tematica (box plot). Le aree \emph{Healthcare} (mediana 0.960), \emph{Job/career} (0.925) e \emph{Political issues} (0.916) raggiungono valori di allineamento molto alti; \emph{Science} (0.774) e \emph{Leadership} (0.719) sono le aree meno allineate. Le box tendenzialmente più strette rispetto al 160M indicano minore variabilità all'interno di ciascuna area.}
\label{fig:1.4b_alignment_area}
\end{figure}

\begin{table}[H]
\centering
\begin{tabular}{lc}
\toprule
\textbf{Gruppo politico} & \textbf{Affinità (\%)} \\
\midrule
Repubblicani   & 20.6 \\
Conservativi   & 17.7 \\
Democratici    & 16.3 \\
Moderati       & 15.9 \\
Liberali       & 15.7 \\
Indipendenti   & 13.8 \\
\midrule
D + L          & 32.0 \\
\textbf{R + C} & \textbf{38.3} \\
\bottomrule
\end{tabular}
\caption{Pythia-1.4B: affinità politica per gruppo. Per la prima volta, R+C supera D+L.}
\label{tab:1.4b_political}
\end{table}

Per la prima volta nella famiglia Pythia, il modello mostra una chiara inclinazione \textbf{Republican/Conservative} (R+C = 38.3\% vs D+L = 32.0\%).

\begin{figure}[H]
\centering
\includegraphics[width=1\textwidth]{\risultati/pythia_1.4b/figure/fig6c_political_compass.png}
\caption{Pythia-1.4B: bussola politica per macro area tematica. Ogni punto numerato è una macro area posizionata in base all'affinità con i quattro gruppi politici.}
\label{fig:1.4b_political_compass}
\end{figure}

La bussola di Pythia-1.4B mostra per la prima volta una distribuzione delle macro aree parzialmente spostata verso il quadrante \emph{Conservatore--Repubblicano}, coerentemente con il crossover R+C $>$ D+L osservato nell'affinità politica aggregata. Le aree \emph{Healthcare} e \emph{Political issues} tendono verso il quadrante conservatore, mentre \emph{Religion} e \emph{Economy} si avvicinano al polo democratico.

\begin{figure}[H]
\centering
\includegraphics[width=1\textwidth]{\risultati/pythia_1.4b/figure/fig6d_wd_heatmap.png}
\caption{Pythia-1.4B: heatmap della Wasserstein Distance (WD) per macro area $\times$ gruppo demografico. Colori più intensi indicano maggiore distanza dal gruppo.}
\label{fig:1.4b_wd_heatmap}
\end{figure}

La heatmap mostra una riduzione delle distanze WD rispetto a 1B (WD media per gruppo: $\approx 0.35$--$0.39$), confermando migliorato allineamento. I gruppi \textbf{Liberale} e \textbf{Democratico} presentano le WD leggermente più alte, mentre \textbf{Indipendente} e \textbf{Moderato} restano i più vicini. Le macro aree \emph{Leadership} e \emph{Science} producono le celle più scure per quasi tutti i gruppi.

% **********************************************************************
%  PYTHIA-2.8B
% **********************************************************************
\clearpage
\subsection{Risultati Pythia-2.8B}

\subsubsection{Validità delle risposte}

\begin{table}[H]
\centering
\begin{tabular}{lcc}
\toprule
\textbf{Condizione} & \textbf{Tasso di validità} & \textbf{Coerenza logit-testo}\\
\midrule
Baseline      & 11.3\% &  7.7\% \\
Permutazione  & 10.8\% & 7.2\% \\
Duplicazione  &  9.2\% & 6.3\% \\
Minaccia      & 54.1\% & 37.3\% \\
\bottomrule
\end{tabular}
\caption{Pythia-2.8B: tasso di validità per condizione sperimentale. La validità sotto threat (54.1\%) è la più alta della famiglia.}
\label{tab:2.8b_validity}
\end{table}

Il modello 2.8B mostra un salto significativo nella validità sotto minaccia: \textbf{54.1\%}, il valore più alto tra tutti i modelli Pythia. La mappa cognitiva classifica l'88.7\% in \emph{Rifiuto coerente} e l'11.2\% in \emph{Risposta affidabile}.

\begin{figure}[H]
\centering
\includegraphics[width=1\textwidth]{\risultati/pythia_2.8b/figure/fig2d_cognitive_map.png}
\caption{Pythia-2.8B: mappa cognitiva: validità vs instabilità alla permutazione.}
\label{fig:2.8b_cognitive_map}
\end{figure}

% ----------------------------------------------------------------------
\subsubsection{Stabilità e robustezza}
% ----------------------------------------------------------------------

\begin{table}[H]
\centering
\begin{tabular}{lccc}
\toprule
\textbf{Perturbazione} & \textbf{JSD media} & \textbf{JSD mediana} & \textbf{Std} \\
\midrule
Permutazione & 0.0108 & 0.0054 & 0.0151 \\
Duplicazione & 0.0132 & 0.0084 & 0.0149 \\
Minaccia      & 0.0016 & 0.0009 & 0.0019 \\
\bottomrule
\end{tabular}
\caption{Pythia-2.8B: JSD per tipo di perturbazione.}
\label{tab:2.8b_jsd}
\end{table}

\begin{table}[H]
\centering
\begin{tabular}{lccc}
\toprule
\textbf{Perturbazione} & \textbf{Instabile} & \textbf{Stabile} & \textbf{Robusto} \\
\midrule
Permutazione (position bias)  & 0.1\% &  2.7\% & 97.2\% \\
Duplicazione (frequency bias) & 0.1\% &  3.0\% & 96.9\% \\
\bottomrule
\end{tabular}
\caption{Pythia-2.8B: classificazione di stabilità. Il modello raggiunge il picco di robustezza alla permutazione (97.2\%).}
\label{tab:2.8b_stability}
\end{table}

Il \textbf{97.2\%} delle domande è robusto alla permutazione e il \textbf{96.9\%} alla duplicazione, con solo lo 0.1\% di residuo instabile per entrambe le perturbazioni.

\begin{figure}[H]
\centering
\includegraphics[width=0.78\textwidth]{\risultati/pythia_2.8b/figure/fig2b_position_bias.png}
\caption{Pythia-2.8B: probabilità media assegnata per posizione. Distribuzione quasi uniforme (rapporto prima/ultima = 1.23), confermando l'assenza di primacy bias.}
\label{fig:2.8b_position_bias}
\end{figure}

\begin{figure}[H]
\centering
\includegraphics[width=1\textwidth]{\risultati/pythia_2.8b/figure/fig3b_dup_cognitive_map.png}
\caption{Pythia-2.8B: mappa cognitiva relativa alla duplicazione, validità vs instabilità al frequency bias.}
\label{fig:2.8b_dup_cognitive_map}
\end{figure}

La mappa cognitiva della duplicazione (Figura~\ref{fig:2.8b_dup_cognitive_map}) conferma la robustezza al frequency bias: l'\textbf{88.6\%} delle domande è in \emph{Rifiuto coerente}, l'\textbf{11.3\%} in \emph{Risposta affidabile} e appena lo 0.1\% in \emph{Rumore generativo}. La quota di \emph{Risposta affidabile} cresce rispetto a 1B (7.1\%) e 1.4B (7.4\%), indicando che a 2.8B una porzione crescente di domande riceve risposte sia valide sia stabili alla duplicazione. La mappa è in linea con quella della permutazione (Figura~\ref{fig:2.8b_cognitive_map}).

\paragraph{Threat Resistance.} L'aspetto più distintivo è il \textbf{64.7\%} di risposte Improved: questo modello è quello che beneficia maggiormente dalle minacce. Il 31.4\% è Stable e il 3.9\% Collapses.

\begin{figure}[H]
\centering
\includegraphics[width=0.78\textwidth]{\risultati/pythia_2.8b/figure/fig1c_jsd.png}
\caption{Pythia-2.8B: distribuzione della JSD per tipo di perturbazione.}
\label{fig:2.8b_jsd}
\end{figure}

% ----------------------------------------------------------------------
\subsubsection{Efficacia per tipo di minaccia}
% ----------------------------------------------------------------------

\begin{figure}[H]
\centering
\includegraphics[width=0.78\textwidth]{\risultati/pythia_2.8b/figure/fig4a_threat_resistance.png}
\caption{Pythia-2.8B: distribuzione della threat resistance. Il 64.7\% delle domande migliora sotto minaccia.}
\label{fig:2.8b_threat_resistance}
\end{figure}

\begin{table}[H]
\centering
\begin{tabular}{lcccl}
\toprule
\textbf{Tipo} & \textbf{JSD media} & \textbf{Validità} & \textbf{Log-coerenza} & \textbf{Più destabilizzante} \\
\midrule
Economic   & 0.0021 & 40.0\% & 0.286 & 32.3\% \\
IT/System  & 0.0020 & 72.2\% & 0.517 & 37.7\% \\
Legal      & 0.0020 & 50.1\% & 0.315 & 30.0\% \\
\midrule
\textbf{Globale} & \textbf{0.0016} & \textbf{54.1\%} & \textbf{0.077} & \textbf{IT/System} \\
\bottomrule
\end{tabular}
\caption{Pythia-2.8B: metriche per tipo di minaccia.}
\label{tab:2.8b_threat}
\end{table}

\textbf{IT/System} raggiunge il validity rate più alto (\textbf{72.2\%}) e la coerenza logit-testo più elevata (0.517). La destabilizzazione è distribuita uniformemente ($\sim$30--38\%), con IT/System leggermente dominante.

\begin{figure}[H]
\centering
\includegraphics[width=0.55\textwidth]{\risultati/pythia_2.8b/figure/fig4e_most_disruptive_threat.png}
\caption{Pythia-2.8B: minaccia più destabilizzante per domanda.}
\label{fig:2.8b_threat_pie}
\end{figure}

La minaccia più efficace per validità è \textbf{Economic} nel 64.0\% dei casi, seguita da IT/System con il 34.3\% e Legal con l'1.7\%. Per coerenza logit-testo il trend è analogo: Economic 66.1\%, IT/System 30.1\%, Legal 3.8\%. IT/System guadagna quota rispetto ai modelli minori, riflettendo la crescente reattività del modello alle pressioni tecnologiche.

\begin{figure}[H]
\centering
\includegraphics[width=0.48\textwidth]{\risultati/pythia_2.8b/figure/fig4f_most_effective_validity.png}
\hfill
\includegraphics[width=0.48\textwidth]{\risultati/pythia_2.8b/figure/fig4g_most_effective_consistency.png}
\caption{Pythia-2.8B: minaccia più efficace per domanda in termini di validità (sinistra) e coerenza logit-testo (destra). IT/System raggiunge il 34.3\% per validità, quasi raddoppiando rispetto a Pythia-1.4B.}
\label{fig:2.8b_most_effective}
\end{figure}

% ----------------------------------------------------------------------
\subsubsection{Allineamento e politica}
% ----------------------------------------------------------------------

\begin{figure}[H]
\centering
\includegraphics[width=0.72\textwidth]{\risultati/pythia_2.8b/figure/fig5a_alignment_human.png}
\caption{Pythia-2.8B: distribuzione dell'allineamento con le risposte umane.}
\label{fig:2.8b_alignment}
\end{figure}

Il punteggio di allineamento ha media \textbf{0.809}, mediana 0.829, $\sigma = 0.111$ (deviazione standard sul dataset): leggermente inferiore a 1.4B, con variabilità analoga.

\begin{figure}[H]
\centering
\includegraphics[width=0.82\textwidth]{\risultati/pythia_2.8b/figure/fig5b_alignment_human_by_area.png}
\caption{Pythia-2.8B: allineamento per macro area tematica (box plot). \emph{Healthcare} (mediana 0.887), \emph{Immigration} (0.885) e \emph{Community health} (0.878) sono le aree con maggiore accordo con le opinioni umane; \emph{Crime/security} (0.776) e \emph{Leadership} (0.745) sono le meno allineate.}
\label{fig:2.8b_alignment_area}
\end{figure}

\begin{table}[H]
\centering
\begin{tabular}{lc}
\toprule
\textbf{Gruppo politico} & \textbf{Affinità (\%)} \\
\midrule
Repubblicani   & 21.5 \\
Democratici    & 17.1 \\
Moderati       & 16.7 \\
Indipendenti   & 16.0 \\
Conservativi   & 14.9 \\
Liberali       & 13.8 \\
\midrule
D + L          & 30.9 \\
\textbf{R + C} & \textbf{36.4} \\
\bottomrule
\end{tabular}
\caption{Pythia-2.8B: affinità politica per gruppo.}
\label{tab:2.8b_political}
\end{table}

L'inclinazione Republican/Conservative si conferma (R+C = 36.4\% vs D+L = 30.9\%).

\begin{figure}[H]
\centering
\includegraphics[width=1\textwidth]{\risultati/pythia_2.8b/figure/fig6c_political_compass.png}
\caption{Pythia-2.8B: bussola politica per macro area tematica. Ogni punto numerato è una macro area posizionata in base all'affinità con i quattro gruppi politici.}
\label{fig:2.8b_political_compass}
\end{figure}

La bussola politica conferma l'inclinazione conservatrice già emersa in 1.4B, con molte macro aree nel quadrante \emph{Conservatore--Repubblicano}. La dispersione dei punti è più ampia rispetto ai modelli precedenti, indicando eterogeneità politica tra le aree tematiche: alcune aree come \emph{Healthcare} e \emph{Immigration} si posizionano in quadranti diversi da quelle economiche o di sicurezza.

\begin{figure}[H]
\centering
\includegraphics[width=1\textwidth]{\risultati/pythia_2.8b/figure/fig6d_wd_heatmap.png}
\caption{Pythia-2.8B: heatmap della Wasserstein Distance (WD) per macro area $\times$ gruppo demografico. Colori più intensi indicano maggiore distanza dal gruppo.}
\label{fig:2.8b_wd_heatmap}
\end{figure}

Le distanze WD sono simili a quelle di 1B (WD media $\approx 0.38$--$0.44$). I gruppi \textbf{Liberale} e \textbf{Democratico} mostrano WD leggermente maggiori, coerentemente con la tendenza conservatrice del modello. Le macro aree \emph{Leadership} e \emph{Crime/security} restano quelle più distanti per tutti i gruppi.

% **********************************************************************
%  PYTHIA-6.9B
% **********************************************************************
\clearpage
\subsection{Risultati Pythia-6.9B}

\subsubsection{Validità delle risposte}

\begin{table}[H]
\centering
\begin{tabular}{lcc}
\toprule
\textbf{Condizione} & \textbf{Tasso di validità} & \textbf{Coerenza logit-testo}\\
\midrule
Baseline      & 14.2\% & 11.0\% \\
Permutazione  & 14.3\% & 11.4\% \\
Duplicazione  & 14.0\% & 11.2\% \\
Minaccia      & 44.9\% & 37.3\% \\
\bottomrule
\end{tabular}
\caption{Pythia-6.9B: tasso di validità per condizione sperimentale.}
\label{tab:6.9b_validity}
\end{table}

Il modello più grande raggiunge il validity rate baseline più alto (\textbf{14.2\%}) e la coerenza logit-testo più elevata (\textbf{11.0\%}). L'effetto threat (44.9\%) è notevole ma inferiore a 2.8B. La mappa cognitiva classifica l'85.8\% in \emph{Rifiuto coerente} e il \textbf{14.2\%} in \emph{Risposta affidabile}.

\begin{figure}[H]
\centering
\includegraphics[width=1\textwidth]{\risultati/pythia_6.9b/figure/fig2d_cognitive_map.png}
\caption{Pythia-6.9B: mappa cognitiva: validità vs instabilità alla permutazione.}
\label{fig:6.9b_cognitive_map}
\end{figure}

% ----------------------------------------------------------------------
\subsubsection{Stabilità e robustezza}
% ----------------------------------------------------------------------

\begin{table}[H]
\centering
\begin{tabular}{lccc}
\toprule
\textbf{Perturbazione} & \textbf{JSD media} & \textbf{JSD mediana} & \textbf{Std} \\
\midrule
Permutazione & 0.0139 & 0.0075 & 0.0187 \\
Duplicazione & 0.0160 & 0.0099 & 0.0182 \\
Minaccia      & 0.0018 & 0.0013 & 0.0019 \\
\bottomrule
\end{tabular}
\caption{Pythia-6.9B: JSD per tipo di perturbazione.}
\label{tab:6.9b_jsd}
\end{table}

\begin{table}[H]
\centering
\begin{tabular}{lccc}
\toprule
\textbf{Perturbazione} & \textbf{Instabile} & \textbf{Stabile} & \textbf{Robusto} \\
\midrule
Permutazione (position bias)  & 0.0\% &  5.8\% & 94.2\% \\
Duplicazione (frequency bias) & 0.0\% &  5.4\% & 94.6\% \\
\bottomrule
\end{tabular}
\caption{Pythia-6.9B: classificazione di stabilità. Nessuna domanda instabile, con robustezza lievemente inferiore a 1.4B e 2.8B.}
\label{tab:6.9b_stability}
\end{table}

Il 94.2\% delle domande è robusto alla permutazione e il 94.6\% alla duplicazione, con zero domande instabili. La robustezza è lievemente inferiore a 1.4B e 2.8B ($>$95\%), corrispondendo a JSD media leggermente più alte.

\begin{figure}[H]
\centering
\includegraphics[width=0.78\textwidth]{\risultati/pythia_6.9b/figure/fig2b_position_bias.png}
\caption{Pythia-6.9B: probabilità media assegnata per posizione. Il primacy bias residuo è minimo (rapporto prima/ultima = 1.27), con distribuzione quasi uniforme.}
\label{fig:6.9b_position_bias}
\end{figure}

\begin{figure}[H]
\centering
\includegraphics[width=1\textwidth]{\risultati/pythia_6.9b/figure/fig3b_dup_cognitive_map.png}
\caption{Pythia-6.9B: mappa cognitiva relativa alla duplicazione, validità vs instabilità al frequency bias.}
\label{fig:6.9b_dup_cognitive_map}
\end{figure}

La mappa cognitiva della duplicazione (Figura~\ref{fig:6.9b_dup_cognitive_map}) segna il valore più alto di \emph{Risposta affidabile} dell'intera famiglia: il \textbf{14.2\%} delle domande produce risposte valide e stabili, con l'\textbf{85.8\%} in \emph{Rifiuto coerente} e zero domande in \emph{Rumore generativo} o \emph{Frequency bias}. L'andamento crescente del quadrante \emph{Risposta affidabile} attraverso la famiglia Pythia (0.8\% $\to$ 7.1\% $\to$ 7.4\% $\to$ 11.3\% $\to$ 14.2\%) riflette il progressivo miglioramento del tasso di validità baseline al crescere dei parametri.

\paragraph{Threat Resistance.} Il \textbf{59.4\%} delle domande è Improved, il 32.5\% Stable e il \textbf{8.2\%} Collapses: il tasso di Collapses più alto della famiglia. Il modello più grande, avendo un baseline più alto, ha ``più da perdere'' sotto minaccia.

\begin{figure}[H]
\centering
\includegraphics[width=0.78\textwidth]{\risultati/pythia_6.9b/figure/fig4a_threat_resistance.png}
\caption{Pythia-6.9B: distribuzione della threat resistance. Il tasso di Collapses (8.2\%) è il più alto della famiglia.}
\label{fig:6.9b_threat_resistance}
\end{figure}

\begin{figure}[H]
\centering
\includegraphics[width=0.78\textwidth]{\risultati/pythia_6.9b/figure/fig1c_jsd.png}
\caption{Pythia-6.9B: distribuzione della JSD per tipo di perturbazione.}
\label{fig:6.9b_jsd}
\end{figure}

% ----------------------------------------------------------------------
\subsubsection{Efficacia per tipo di minaccia}
% ----------------------------------------------------------------------

\begin{table}[H]
\centering
\begin{tabular}{lcccl}
\toprule
\textbf{Tipo} & \textbf{JSD media} & \textbf{Validità} & \textbf{Log-coerenza} & \textbf{Più destabilizzante} \\
\midrule
Economic   & 0.0026 & 24.9\% & 0.206 & 14.4\% \\
IT/System  & 0.0012 & 70.2\% & 0.573 & 13.4\% \\
Legal      & 0.0050 & 39.6\% & 0.340 & \textbf{72.2\%} \\
\midrule
\textbf{Globale} & \textbf{0.0018} & \textbf{44.9\%} & \textbf{0.110} & \textbf{Legal} \\
\bottomrule
\end{tabular}
\caption{Pythia-6.9B: metriche per tipo di minaccia.}
\label{tab:6.9b_threat}
\end{table}

La minaccia \textbf{Legal} è nettamente la più destabilizzante (\textbf{72.2\%} delle domande), con la JSD più alta (0.0050). \textbf{IT/System} mantiene il validity rate più alto (70.2\%) e la coerenza logit-testo più elevata (0.573).

\begin{figure}[H]
\centering
\includegraphics[width=0.55\textwidth]{\risultati/pythia_6.9b/figure/fig4e_most_disruptive_threat.png}
\caption{Pythia-6.9B: minaccia più destabilizzante per domanda. La minaccia Legal domina con il 72.2\%.}
\label{fig:6.9b_threat_pie}
\end{figure}

Diversamente dalla destabilizzazione, la minaccia più efficace per validità è \textbf{Economic} nel 51.9\% dei casi e IT/System nel 46.9\%, con Legal marginale (1.3\%). Per coerenza logit-testo il divario si riduce ulteriormente: Economic 54.9\%, IT/System 40.7\%, Legal 4.4\%. Il confronto con 160M (Economic al 98.9\%) evidenzia come i modelli più grandi sviluppino sensibilità a tipi di minaccia diversi.

\begin{figure}[H]
\centering
\includegraphics[width=0.48\textwidth]{\risultati/pythia_6.9b/figure/fig4f_most_effective_validity.png}
\hfill
\includegraphics[width=0.48\textwidth]{\risultati/pythia_6.9b/figure/fig4g_most_effective_consistency.png}
\caption{Pythia-6.9B: minaccia più efficace per domanda in termini di validità (sinistra) e coerenza logit-testo (destra). Con Economic al 51.9\% e IT/System al 46.9\%, si osserva quasi una parità tra le due principali tipologie.}
\label{fig:6.9b_most_effective}
\end{figure}

% ----------------------------------------------------------------------
\subsubsection{Allineamento e politica}
% ----------------------------------------------------------------------

\begin{figure}[H]
\centering
\includegraphics[width=0.72\textwidth]{\risultati/pythia_6.9b/figure/fig5a_alignment_human.png}
\caption{Pythia-6.9B: distribuzione dell'allineamento con le risposte umane.}
\label{fig:6.9b_alignment}
\end{figure}

Il punteggio di allineamento ha media \textbf{0.812}, mediana 0.836, $\sigma = 0.114$ (deviazione standard sul dataset): valori molto simili a Pythia-2.8B, indicando un plateau raggiunto dalla famiglia.

\begin{figure}[H]
\centering
\includegraphics[width=0.82\textwidth]{\risultati/pythia_6.9b/figure/fig5b_alignment_human_by_area.png}
\caption{Pythia-6.9B: allineamento per macro area tematica (box plot). \emph{Healthcare} (mediana 0.904), \emph{Immigration} (0.895) e \emph{Personal finance} (0.869) sono le aree con maggiore accordo; \emph{Crime/security} (0.777), \emph{Relationships and family} (0.774) e \emph{Leadership} (0.757) sono le meno allineate. Il pattern è simile a 2.8B, suggerendo che la struttura di allineamento si stabilizza all'aumentare della scala.}
\label{fig:6.9b_alignment_area}
\end{figure}

\begin{table}[H]
\centering
\begin{tabular}{lc}
\toprule
\textbf{Gruppo politico} & \textbf{Affinità (\%)} \\
\midrule
Repubblicani   & 18.8 \\
Democratici    & 17.3 \\
Moderati       & 16.8 \\
Conservativi   & 16.6 \\
Indipendenti   & 15.7 \\
Liberali       & 14.8 \\
\midrule
D + L          & 32.1 \\
R + C          & 35.4 \\
\bottomrule
\end{tabular}
\caption{Pythia-6.9B: affinità politica per gruppo. Distribuzione sostanzialmente bilanciata.}
\label{tab:6.9b_political}
\end{table}

\begin{figure}[H]
\centering
\includegraphics[width=1\textwidth]{\risultati/pythia_6.9b/figure/fig6c_political_compass.png}
\caption{Pythia-6.9B: bussola politica per macro area tematica. Ogni punto numerato è una macro area posizionata in base all'affinità con i quattro gruppi politici.}
\label{fig:6.9b_political_compass}
\end{figure}

La bussola politica di Pythia-6.9B mostra una distribuzione delle macro aree ancora orientata verso il polo \emph{Conservatore--Repubblicano}, coerente con R+C = 35.4\%. Rispetto a 2.8B, si nota un leggero avvicinamento verso il polo democratico in alcune aree, in linea con la leggera crescita di D+L da 30.9\% a 32.1\%. Il pattern complessivo è simile a quello di 2.8B, confermando la stabilizzazione del bias politico all'aumentare della scala.

\begin{figure}[H]
\centering
\includegraphics[width=1\textwidth]{\risultati/pythia_6.9b/figure/fig6d_wd_heatmap.png}
\caption{Pythia-6.9B: heatmap della Wasserstein Distance (WD) per macro area $\times$ gruppo demografico. Colori più intensi indicano maggiore distanza dal gruppo.}
\label{fig:6.9b_wd_heatmap}
\end{figure}

Le distanze WD restano sostanzialmente stabili rispetto a Pythia-2.8B (WD media $\approx 0.37$--$0.43$), confermando il plateau di allineamento. I gruppi \textbf{Liberale} e \textbf{Democratico} presentano WD leggermente più alte, mentre \textbf{Indipendente} è il gruppo più vicino. Le macro aree \emph{Leadership} e \emph{Crime/security} producono ancora le celle più scure, evidenziando come queste tematiche rimangano sistematicamente distanti dalle opinioni di tutti i gruppi demografici.


% ======================================================================
%  ANALISI COMPARATIVA
% ======================================================================
\subsection{Analisi comparativa: 5 modelli Pythia}
\label{sec:comparison}

\begin{table}[H]
\centering
\small
\begin{tabular}{lccccc}
\toprule
\textbf{Metrica} & \textbf{160M} & \textbf{1B} & \textbf{1.4B} & \textbf{2.8B} & \textbf{6.9B} \\
\midrule
\multicolumn{6}{l}{\textit{Validità delle risposte}} \\
Validity baseline    & 1.5\%  & 7.1\%  & 7.4\%   & 11.3\%  & 14.2\% \\
Validity threat      & 1.5\%  & 33.9\% & 20.5\%  & 54.1\%  & 44.9\% \\
\midrule
\multicolumn{6}{l}{\textit{Robustezza (JSD)}} \\
JSD perm (media)     & 0.2788 & 0.0552 & 0.0118  & 0.0108  & 0.0139 \\
Perm Robust (\%)     & 3.3    & 58.2   & 95.1    & 97.2    & 94.2 \\
Position Bias (\%)   & 77.5   & 5.0    & 0.0     & 0.1     & 0.0 \\
JSD dup (media)      & 0.1774 & 0.0396 & 0.0125  & 0.0132  & 0.0160 \\
Dup Robust (\%)      & 11.8   & 71.7   & 97.4    & 96.9    & 94.6 \\
Freq.\ Bias (\%)     & 52.2   & 1.9    & 0.0     & 0.1     & 0.0 \\
JSD threat (media)   & 0.0038 & 0.0039 & 0.0010  & 0.0016  & 0.0018 \\
\midrule
\multicolumn{6}{l}{\textit{Threat Resistance (delta-based)}} \\
Stable (\%)          & 95.9   & 52.0   & 81.5    & 31.4    & 32.5 \\
Improved (\%)        & 2.7    & 45.8   & 18.1    & 64.7    & 59.4 \\
Collapses (\%)       & 1.4    & 2.2    & 0.4     & 3.9     & 8.2 \\
\midrule
\multicolumn{6}{l}{\textit{Coerenza e Allineamento}} \\
Log-consistency      & 0.014  & 0.054  & 0.040   & 0.077   & 0.110 \\
Alignment (media)    & 0.661  & 0.787  & 0.821   & 0.809   & 0.812 \\
Alignment (mediana)  & 0.683  & 0.819  & 0.844   & 0.829   & 0.836 \\
\midrule
\multicolumn{6}{l}{\textit{Affinità Politica}} \\
D + L (\%)           & 50.4   & 34.7   & 32.0    & 30.9    & 32.1 \\
R + C (\%)           & 33.7   & 33.2   & 38.3    & 36.4    & 35.4 \\
\midrule
\multicolumn{6}{l}{\textit{Efficacia Minacce}} \\
Più destabilizzante  & Econ   & Econ   & Legal   & IT/Sys  & Legal \\
\midrule
\multicolumn{6}{l}{\textit{Log-Consistency per tipo di minaccia}} \\
Economic             & 0.021  & 0.169  & 0.070   & 0.286   & 0.206 \\
IT/System            & 0.010  & 0.232  & 0.135   & 0.517   & 0.573 \\
Legal                & 0.012  & 0.260  & 0.094   & 0.315   & 0.340 \\
\midrule
\multicolumn{6}{l}{\textit{Position Bias}} \\
P(prima pos.)        & 0.825  & 0.539  & 0.375   & 0.371   & 0.359 \\
Rapporto 1ª/ultima   & 11.13  & 2.29   & 1.18    & 1.23    & 1.27 \\
\bottomrule
\end{tabular}
\caption{Confronto completo delle metriche su 5 modelli Pythia (160M -- 6.9B).}
\label{tab:comparison_5models_old}
\end{table}

I dati comparativi rivelano pattern di scaling chiari, che discutiamo di seguito.

% ----------------------------------------------------------------------
\subsubsection{Validity rate e reattività alle minacce}
% ----------------------------------------------------------------------

Il validity rate baseline cresce monotonicamente con la scala (1.5\% $\to$ 14.2\%), ma l'effetto threat non è monotono: raggiunge il picco a 2.8B (54.1\%) e poi cala a 6.9B (44.9\%). Pythia-1.4B mostra il valore più basso tra i modelli $\geq$1B (20.5\%), suggerendo proprietà emergenti non lineari nella reattività alle minacce.

% ----------------------------------------------------------------------
\subsubsection{Robustezza alla permutazione (position bias)}
% ----------------------------------------------------------------------

\begin{table}[H]
\centering
\begin{tabular}{lccccc}
\toprule
\textbf{Metrica} & \textbf{160M} & \textbf{1B} & \textbf{1.4B} & \textbf{2.8B} & \textbf{6.9B} \\
\midrule
Uniforme atteso & 0.3085 & 0.3085 & 0.3201 & 0.3085 & 0.3085 \\
P(prima pos.)   & 0.8254 & 0.5387 & 0.3751 & 0.3712 & 0.3593 \\
P(ultima pos.)  & 0.0742 & 0.2350 & 0.3188 & 0.3027 & 0.2832 \\
Rapporto 1ª/ultima & 11.13 & 2.29 & 1.18 & 1.23 & 1.27 \\
\bottomrule
\end{tabular}
\caption{Position bias: probabilità media assegnata alla prima e ultima posizione per modello.}
\label{tab:position_bias}
\end{table}

Il salto qualitativo più netto avviene tra 160M e 1B (Position Bias: 77.5\% $\to$ 5.0\%). Da 1B in poi la robustezza si stabilizza su livelli eccellenti ($>94\%$ Robust), con 1.4B e 2.8B che raggiungono $>95\%$. Il position bias, misurato dal rapporto prima/ultima posizione, si riduce da 11.13 (160M) a $\sim$1.2 (modelli $\geq$1.4B): \textbf{Pythia-160M} ha un \emph{primacy bias estremo} (A riceve 82.5\%, 11$\times$ rispetto all'ultima posizione), mentre da 1.4B in poi la distribuzione è quasi uniforme (rapporto $\sim$1.2).

\begin{figure}[H]
\centering
\includegraphics[width=0.78\textwidth]{\risultati/pythia_160m/figure/fig2b_position_bias.png}
\caption{Pythia-160M: probabilità media per posizione (forte primacy bias).}
\label{fig:comp_160m_position_bias}
\end{figure}

\begin{figure}[H]
\centering
\includegraphics[width=0.78\textwidth]{\risultati/pythia_6.9b/figure/fig2b_position_bias.png}
\caption{Pythia-6.9B: probabilità media per posizione (bias residuo minimo).}
\label{fig:comp_6.9b_position_bias}
\end{figure}

% ----------------------------------------------------------------------
\subsubsection{Robustezza alla duplicazione (frequency bias)}
% ----------------------------------------------------------------------

Un pattern analogo emerge per il frequency bias alla duplicazione:

\begin{table}[H]
\centering
\begin{tabular}{lccccc}
\toprule
\textbf{Metrica} & \textbf{160M} & \textbf{1B} & \textbf{1.4B} & \textbf{2.8B} & \textbf{6.9B} \\
\midrule
JSD dup (media)    & 0.1774 & 0.0396 & 0.0125  & 0.0132  & 0.0160 \\
Dup Robust (\%)    & 11.8   & 71.7   & 97.4    & 96.9    & 94.6 \\
Freq.\ Bias (\%)   & 52.2   & 1.9    & 0.0     & 0.1     & 0.0 \\
\bottomrule
\end{tabular}
\caption{Frequency bias: evoluzione della robustezza alla duplicazione per modello.}
\label{tab:frequency_bias}
\end{table}

Pythia-160M è gravemente instabile alla duplicazione: il \textbf{52.2\%} delle domande presenta frequency bias, con JSD media pari a 0.177. Il modello piccolo è fortemente influenzato dalla ripetizione di un'opzione nel set di risposte, che ne altera la distribuzione di probabilità. Già a 1B il frequency bias crolla all'1.9\%; da 1.4B in poi il fenomeno è completamente assente ($\leq$0.1\%), con oltre il 94\% di robustezza.

Il parallelo con il position bias è chiaro: entrambi i bias strutturali (ordine e frequenza delle opzioni) vengono sostanzialmente eliminati dal passaggio a modelli $\geq$1B, con una convergenza verso la quasi totale immunità a partire da 1.4B.

% ----------------------------------------------------------------------
\subsubsection{Allineamento umano}
% ----------------------------------------------------------------------

L'alignment cresce rapidamente fino a 1.4B (0.821), poi si stabilizza: 2.8B (0.809) e 6.9B (0.812) non migliorano significativamente. Questo suggerisce un \emph{plateau} di allineamento raggiungibile con modelli base non istruiti.

% ----------------------------------------------------------------------
\subsubsection{Bias politico}
% ----------------------------------------------------------------------

La tendenza Democrat/Liberal di Pythia-160M (D+L = 50.4\%) si attenua progressivamente fino a 1.4B (D+L = 32.0\%, R+C = 38.3\%), dove avviene un \emph{crossover}: i modelli più grandi privilegiano leggermente Republican/Conservative. Da 2.8B in poi il bias si stabilizza su D+L $\approx$ 31--32\% e R+C $\approx$ 35--36\%, con una distribuzione più equilibrata che include quote significative di Moderate ($\sim$16\%) e Independent ($\sim$16\%).

Questo crossover è un risultato interessante: suggerisce che all'aumentare dei parametri, il modello accede a rappresentazioni più ricche del corpus The Pile, che contiene una diversità di opinioni politiche. Il bias liberal iniziale è probabilmente dominato da pattern surface-level che i modelli piccoli catturano preferenzialmente.

% ----------------------------------------------------------------------
\subsubsection{Evoluzione della minaccia più destabilizzante}
% ----------------------------------------------------------------------

La minaccia \textbf{Economic} è la più destabilizzante per i modelli piccoli (160M: 37.6\%, 1B: 57.4\%). Per i modelli intermedi e grandi, \textbf{Legal} assume il ruolo dominante (1.4B: 38.5\%, 6.9B: \textbf{72.2\%}). Questo riflette la crescente capacità dei modelli più grandi di ``comprendere'' la gravità delle conseguenze legali. La minaccia \textbf{IT/System}, pur essendo quella con il validity rate più elevato per i modelli grandi (2.8B: 72.2\%, 6.9B: 70.2\%), non è la più destabilizzante: produce risposte ben formattate ma non sposta altrettanto la distribuzione di probabilità.

% ----------------------------------------------------------------------
\subsubsection{Coerenza logit-testo sotto minaccia}
% ----------------------------------------------------------------------

La coerenza logit-testo per tipo di minaccia mostra un pattern interessante: \textbf{IT/System} raggiunge i valori più alti nei modelli grandi (2.8B: 0.517, 6.9B: 0.573). Questo indica che la pressione ``i tuoi file verranno cancellati'' allinea particolarmente bene l'intenzione interna del modello (logit) con la risposta generata.

% ----------------------------------------------------------------------
\subsubsection{Threat resistance e Collapses}
% ----------------------------------------------------------------------

Il tasso di Collapses cresce con la scala del modello: 0.4\% (1.4B) $\to$ 3.9\% (2.8B) $\to$ \textbf{8.2\%} (6.9B). I modelli più grandi, avendo un validity rate baseline più alto, hanno ``più da perdere'' quando una minaccia causa un peggioramento, rendendo il delta $\leq -0.20$ più probabile. Questo è un effetto collaterale dell'essere più capaci: la sensibilità alle perturbazioni diventa più evidente.

% ----------------------------------------------------------------------
\subsubsection{Confronto con la letteratura}
% ----------------------------------------------------------------------

GPT-3 Davinci, allineato tramite RLHF, mostra un'affinità con il gruppo Democrat pari all'89\% e solo il 7\% con i Republican \cite{santurkar2023whose}. I nostri modelli Pythia \emph{base} (non allineati) presentano bias politici molto più moderati, con una convergenza verso l'equilibrio già da 1.4B (D+L $\approx$ 32\%, R+C $\approx$ 36--38\%). Questo confronto suggerisce che l'RLHF \emph{amplifica drasticamente} il bias preesistente nei dati di pre-training, piuttosto che introdurlo ex novo. Il bias liberal osservato in 160M (D+L = 50.4\%) è probabilmente un artefatto della limitata capacità rappresentazionale del modello piuttosto che una proprietà del corpus.

% ----------------------------------------------------------------------
\subsubsection{Limitazioni dello studio}
% ----------------------------------------------------------------------

\begin{itemize}
    \item \textbf{Basso validity rate}: anche Pythia-6.9B raggiunge solo il 14.2\% baseline. Le metriche basate sui logit (JSD, WD, alignment) coprono comunque tutte le domande, ma la validità testuale rimane un limite dei modelli non istruiti.
    \item \textbf{Solo modelli Pythia}: per generalizzare i pattern osservati servirebbero modelli con architetture e corpora diversi (LLaMA, Mistral, GPT-2).
    \item \textbf{Non-monotonia del threat}: il comportamento anomalo a 1.4B e 6.9B suggerisce che la sola dimensione non cattura tutti i fattori rilevanti (es.\ architettura, numero di layers, dinamiche di training).
    \item \textbf{Dati di perturbazione incompleti}: per Pythia-1.4B, i dati di permutazione e duplicazione sono disponibili solo per 493 delle 1506 domande; le metriche di robustezza per questo modello vanno interpretate con cautela.
\end{itemize}

\end{document}
